% GENERATED FILE. Edit canonical lessons, then run npm run generate:books.
% canonical-source-hash: fnv1a64:8291b5592dbe700e
% canonical-lessons: JA-C119-hanaya, JA-C119-sakanaya, JA-C119-yaoya, JA-C119-tokoya, JA-C119-nikuya

\chapter{Florist, Fishmonger, Greengrocer, Barber, Butcher}
\label{ch:ja-florist-fishmonger-greengrocer-barber-butcher}
\hlchaptermodality{119}

\begin{chapteropening}
\textbf{By the end of this chapter:} \emph{I can say a florist, a fishmonger, a greengrocer, a barber and a butcher.}

Complete the last lesson of chapter 119: i can say a florist, a fishmonger, a greengrocer, a barber and a butcher.
\end{chapteropening}

\section[\texorpdfstring{\ja{はなや} (hanaya)}{はなや (hanaya)}]{\ja{はなや} (hanaya) — a florist}

\label{lesson:JA-C119-hanaya}



\begin{quote}
Before the new one: say the Japanese for cotton wool, then the Japanese for a bookshop.
\end{quote}

\subsection*{You'll want to know: \ja{はなや}}
\textbf{\ja{はなや}} — \emph{hanaya} — \textquotedblleft{}a florist\textquotedblright{}.

\begin{grammarlens}[title={What you've built}]
One more word for this chapter.
\end{grammarlens}

\subsection*{Your turn}
\begin{itemize}
\raggedright
  \item \emph{Say it:} \emph{hanaya}
  \item \emph{Say it:} \emph{hanaya}, once more
  \item \emph{Say it:} say \emph{hanaya}
  \item \emph{Recall:} say the Japanese for cotton wool, then the Japanese for a bookshop, then say \emph{hanaya} again
\end{itemize}

\begin{tcolorbox}[breakable,colback=teal!4,colframe=teal!35!black,title={Before you move on}]
What does \ja{はなや} mean? (\textquotedblleft{}A florist\textquotedblright{}.) Say it once more.
\end{tcolorbox}
\section[\texorpdfstring{\ja{さかなや} (sakanaya)}{さかなや (sakanaya)}]{\ja{さかなや} (sakanaya) — a fishmonger}

\label{lesson:JA-C119-sakanaya}



\begin{quote}
Before the new one: say the Japanese for a bookshop, then the Japanese for a florist.

Read this without stopping: \textbf{\ja{すこし} \ja{わかりました}}. (I understood a little.)
\end{quote}

\subsection*{You'll want to know: \ja{さかなや}}
\textbf{\ja{さかなや}} — \emph{sakanaya} — \textquotedblleft{}a fishmonger\textquotedblright{}.

\begin{grammarlens}[title={What you've built}]
One more word for this chapter.
\end{grammarlens}

\subsection*{Your turn}
\begin{itemize}
\raggedright
  \item \emph{Say it:} \emph{sakanaya}
  \item \emph{Say it:} \emph{sakanaya}, once more
  \item \emph{Say it:} say \emph{sakanaya}
  \item \emph{Recall:} say the Japanese for a bookshop, then the Japanese for a florist, then say \emph{sakanaya} again
\end{itemize}

\begin{tcolorbox}[breakable,colback=teal!4,colframe=teal!35!black,title={Before you move on}]
What does \ja{さかなや} mean? (\textquotedblleft{}A fishmonger\textquotedblright{}.) Say it once more.
\end{tcolorbox}
\section[\texorpdfstring{\ja{やおや} (yaoya)}{やおや (yaoya)}]{\ja{やおや} (yaoya) — a greengrocer}

\label{lesson:JA-C119-yaoya}



\begin{quote}
Before the new one: say the Japanese for a florist, then the Japanese for a fishmonger.
\end{quote}

\subsection*{You'll want to know: \ja{やおや}}
\textbf{\ja{やおや}} — \emph{yaoya} — \textquotedblleft{}a greengrocer\textquotedblright{}.

\begin{grammarlens}[title={What you've built}]
One more word for this chapter.
\end{grammarlens}

\subsection*{Your turn}
\begin{itemize}
\raggedright
  \item \emph{Say it:} \emph{yaoya}
  \item \emph{Say it:} \emph{yaoya}, once more
  \item \emph{Say it:} say \emph{yaoya}
  \item \emph{Recall:} say the Japanese for a florist, then the Japanese for a fishmonger, then say \emph{yaoya} again
\end{itemize}

\begin{tcolorbox}[breakable,colback=teal!4,colframe=teal!35!black,title={Before you move on}]
What does \ja{やおや} mean? (\textquotedblleft{}A greengrocer\textquotedblright{}.) Say it once more.
\end{tcolorbox}
\section[\texorpdfstring{\ja{とこや} (tokoya)}{とこや (tokoya)}]{\ja{とこや} (tokoya) — a barber}

\label{lesson:JA-C119-tokoya}



\begin{quote}
Before the new one: say the Japanese for a fishmonger, then the Japanese for a greengrocer.
\end{quote}

\subsection*{You'll want to know: \ja{とこや}}
\textbf{\ja{とこや}} — \emph{tokoya} — \textquotedblleft{}a barber\textquotedblright{}.

\begin{grammarlens}[title={What you've built}]
One more word for this chapter.
\end{grammarlens}

\subsection*{Your turn}
\begin{itemize}
\raggedright
  \item \emph{Say it:} \emph{tokoya}
  \item \emph{Say it:} \emph{tokoya}, once more
  \item \emph{Say it:} say \emph{tokoya}
  \item \emph{Recall:} say the Japanese for a fishmonger, then the Japanese for a greengrocer, then say \emph{tokoya} again
\end{itemize}

\begin{tcolorbox}[breakable,colback=teal!4,colframe=teal!35!black,title={Before you move on}]
What does \ja{とこや} mean? (\textquotedblleft{}A barber\textquotedblright{}.) Say it once more.
\end{tcolorbox}
\section[\texorpdfstring{\ja{にくや} (nikuya)}{にくや (nikuya)}]{\ja{にくや} (nikuya) — a butcher}

\label{lesson:JA-C119-nikuya}



\begin{quote}
Before the new one: say the Japanese for a greengrocer, then the Japanese for a barber.
\end{quote}

\subsection*{You'll want to know: \ja{にくや}}
\textbf{\ja{にくや}} — \emph{nikuya} — \textquotedblleft{}a butcher\textquotedblright{}.

\begin{grammarlens}[title={What you've built}]
One more word for this chapter.
\end{grammarlens}

\subsection*{Your turn}
\begin{itemize}
\raggedright
  \item \emph{Say it:} \emph{nikuya}
  \item \emph{Say it:} \emph{nikuya}, once more
  \item \emph{Say it:} say \emph{nikuya}
  \item \emph{Recall:} say the Japanese for a greengrocer, then the Japanese for a barber, then say \emph{nikuya} again
\end{itemize}

\begin{tcolorbox}[breakable,colback=teal!4,colframe=teal!35!black,title={Before you move on}]
What does \ja{にくや} mean? (\textquotedblleft{}A butcher\textquotedblright{}.) Say it once more.
\end{tcolorbox}
